\documentclass[11pt]{article}
\usepackage[margin=1in]{geometry}
\usepackage{booktabs}
\usepackage{url}
\title{Can a closed-loop terminal set be computed in real time?}
\date{Prepared October 4, 2026}
\begin{document}
\maketitle

\section{Question}
The user proposed a terminal set built from the target contract (constant
tangential acceleration $A$ and sideslip $\beta$): the states from which the
nominal path guidance alone never collides with the target. This study asks
whether such a set can be imposed at every 50-ms sample. The controller is
unchanged. The study script is \path{scripts/studyClosedLoopTerminalFeasibility.m}.

\paragraph{Setup.} The nominal guidance is \texttt{nonlinearBicycleModel.nominalFeedback}, the
law that already extends every anchor. For every fifth recorded ego state of the
14 exact-observation runs of the separating-terminal campaign (commit 4b63978),
the ego is rolled out under that guidance and the target under its model until
one of two things happens. Either the rectangles touch, or a model-based
certificate shows that no later collision occurs. The certificate requires the
ego to have converged to the path. With $d_s=R_E+R_C+0.2$~m
(circumscribed radii), it then requires one of the following:
\begin{itemize}
\item straight road, $\beta\neq0$: the target stays on a fixed circle (centre
  $c$, radius $l_r/|\sin\beta|$); the ego is past $c$ along the road and at
  least $d_s$ outside the circle;
\item straight road, $\beta=0$: the vehicles separate, are at least $d_s$
  apart, and the sign of $\Delta^\top Ae$ keeps the separating rate
  nondecreasing until the target stops;
\item curved road: the ego's path circle and the target's whole future path
  are at least $d_s$ apart, or the target is outside the road circle by $d_s$
  and moving outward.
\end{itemize}
The rollout limit is 4000 steps (200~s).

\section{Results}
\paragraph{Tail length (straight road).} Every non-colliding sampled state
reached a certificate. The tail needed 13--99 steps (median per run) and at
most 187 steps (9.4~s). If the free part must first reach the next state whose
rollout is certified, the total horizon is at most 224 nodes. In the
brakingLead runs about half of the states collide under the nominal guidance:
cruising behind a slower lead is not in the set until the ego has overtaken
(\path{CLOSED_LOOP_TERMINAL_FEASIBILITY_20261004/tail-summary.csv}).

\paragraph{Curved road: the set is empty.} In curvedHeadOn and curvedCrossing,
every sampled state, including those long after the encounter, collides under
the model. The re-collision comes 1000--3460 steps (50--173~s) later. The ego
circles the 200-m-radius road, and the target either circles its own
constant-curvature path or crosses the road circle again. Taken over infinite
time, the contract therefore predicts a later encounter, and no state is safe
forever. Any usable terminal set on a curved road needs a finite validity time
for the target model, or the perception premise that a target beyond the 50-m
radius poses no risk.

\paragraph{Cost per tail node.} The guidance costs 17~$\mu$s, one RK4 step
9~$\mu$s and the hold linearization 12~$\mu$s. The guidance Jacobian by
central differences costs 269~$\mu$s, so a 190-node tail costs about 58~ms. The
Jacobian dominates; an analytic or complex-step Jacobian would remove most of
it. The nonlinear rollouts alone took at most 8.3~ms (p95) on the straight road.

\paragraph{Solve time versus horizon (present formulation).}
\begin{center}
\begin{tabular}{rrrrr}
\toprule
$N$ & startup (ms) & conic (ms) & shifted (ms) & conic (ms)\\
\midrule
64 & 86 & 9 & 46 & 12\\
128 & 74 & 18 & 73 & 21\\
256 & 178 & 35 & 171 & 43\\
512 & 513 & 69 & 504 & 86\\
1024 & 1715 & 138 & 1640 & 168\\
\bottomrule
\end{tabular}
\end{center}
The conic time grows linearly, at about 0.17~ms per node. The total grows
faster than linearly, because the MATLAB formulation assigns sparse entries
inside node loops. If the tail nodes were carried as variables, 224 nodes would
cost about 40~ms of conic time alone.

\paragraph{Condensing the tail.} Along the tail the inputs are a fixed
feedback, so every tail state is an affine function of the endpoint state:
$\delta x_i=\Phi_i\,\delta x_M$. Each tail collision or road row then becomes a
row on the six endpoint variables. The set is a polyhedron in $\delta x_M$, and
the conic problem keeps the size of the free horizon. The closed-loop
sensitivities stay bounded. Starting 2--3~m off the path, the transverse part
of $\|\Phi_i\|$ peaks at 2.8--3.6 at 8~m/s and 5.8--9.6 at 15~m/s, then decays
to 0.01--0.04 after 200 steps (\path{phiStudy.m.txt}). Longitudinal position
sensitivity to speed grows linearly, because path phase is free.

\section{Assessment}
Real-time use appears possible on the straight road if the tail is condensed
into rows on the endpoint state and the guidance Jacobian is computed
analytically. The free horizon then stays at its present size. The added work
is a closed-loop rollout (below 10~ms) and about 12 rows per tail node on six
variables. This is an estimate from the measured parts; the condensed problem
was not built or timed. Three issues remain.
\begin{enumerate}
\item On the curved road the infinite-time set is empty under the contract. A
  finite model-validity time, or the perception-radius premise, must be chosen.
\item The terminal set excludes following a slower lead; the controller must
  overtake within the horizon, or the nominal behaviour must include speed
  adaptation.
\item \texttt{nominalFeedback} contains clips and a sign switch, so its
  Jacobian is piecewise. The peak sensitivity of 9.6 at 15~m/s amplifies
  endpoint linearization error along the tail.
\end{enumerate}

\section{Addendum: free steps needed with a closed-loop terminal set}
The user asked whether the terminal-set design can reduce the number of
prediction steps. If the closed-loop continuation is condensed onto the
endpoint state, it adds rows but no decision steps. The free horizon then only
has to reach a state from which the continuation is safe.
\path{scripts/studyTerminalHorizonReduction.m} measures this along the same
recorded trajectories, at every frame before recovery. It finds the smallest
$j\ge k$ such that, from recorded state $j$, a guidance policy keeps at least
the 0.2-m planning clearance, keeps every rectangle corner on the road, and
reaches the no-collision certificate. $j-k$ is the free horizon that would have
sufficed. Two families are compared: return to the path, and return to the
path or to a parallel lane centre ($+3.66$, $-3.66$ and $+7.32$~m). On the
curved road, where the infinite-time set is empty, safety is taken over a
10-s window instead. The recorded states are a proxy for the planned ones.

\begin{center}
\begin{tabular}{rlrrr}
\toprule
speed & scenario & separating (used) & path only & path or lanes\\
\midrule
8 & headOn & 35 & 13 & 0\\
8 & acceleratingHeadOn & 34 & 12 & 0\\
8 & brakingLead & 98 & 148 & 8\\
8 & crossing & 29 & 27 & 27\\
8 & turningCrossing & 80 & 27 & 27\\
8 & curvedHeadOn & 38 & 14 & 0\\
8 & curvedCrossing & 111 & 26 & 26\\
15 & headOn & 23 & 15 & 14\\
15 & acceleratingHeadOn & 22 & 14 & 14\\
15 & brakingLead & 101 & 135 & 26\\
15 & crossing & 31 & 26 & 25\\
15 & turningCrossing & 61 & 31 & 31\\
15 & curvedHeadOn & 49 & 205 & 32\\
15 & curvedCrossing & 31 & 26 & 25\\
\bottomrule
\end{tabular}
\end{center}
The table gives maxima over the encounter frames. The separating column is the
horizon the controller actually used, whose median is its minimum (8 or 16).
The median free horizon is 0 in every run for the path-or-lanes family: most
recorded states are already in that set. The worst case falls from 111 to 32
steps. Following a braking lead needs the lane policies: returning to the path
needs more steps than the present terminal (148 and 135), because the ego is
returned behind the lead. Crossing targets keep 25--31 steps with either
family, because a parallel lane does not avoid a target that crosses all
lanes; policies with other speeds are the untested next lever. The 205-step
path-only maximum in the 15-m/s curved head-on run was not investigated.

An earlier pass of this study only rejected actual contact and did not check
the road; it reported zero free steps more often. Its numbers are superseded
(\path{horizon-reduction.csv}, \path{horizon.log}).

A union of continuation policies is not convex. It would be solved per policy
(five small conic problems with condensed tails) or by choosing the policy on
the anchor. Preferring the path policy whenever it qualifies keeps the terminal
set aligned with the CLF.

\section{Addendum 2: the CLF's own law as the continuation}
The user clarified the proposal. Once the ego is in the terminal set, the CLF
alone must avoid the target, whose motion is known. The fixed guidance is then
the law that defines the CLF: the discrete LQR feedback
$u=u_{\rm ref}+Ke(x)$, whose Riccati matrix is $P$ in $V=e^\top Pe$
(\texttt{nominalReference.gain}). Braking is clipped to its bounds. The study
was repeated with this law (\texttt{Policy="clfLqr"}). It also counted
continuation steps that violate the requested decrease
$V^+-V\le-\eta\,e^\top Qe$ ($\eta=0.5$), and continuations that leave the model
domain. A continuation that leaves the model domain does not qualify.

\begin{center}
\begin{tabular}{rlrrrrr}
\toprule
 & & separating & \multicolumn{2}{c}{CLF law} & decrease & CLF law, $V\le1$\\
speed & scenario & max & med & max & violated & med / max\\
\midrule
8 & headOn & 35 & 0 & 15 & 6.3\% & 85 / 178\\
8 & acceleratingHeadOn & 34 & 0 & 14 & 7.7\% & 86 / 179\\
8 & brakingLead & 98 & 36 & 148 & 3.4\% & 105 / 217\\
8 & crossing & 29 & 0 & 30 & 3.6\% & 85 / 177\\
8 & turningCrossing & 80 & 0 & 33 & 1.4\% & 68 / 161\\
8 & curvedHeadOn & 38 & 0 & 16 & 1.4\% & 86 / 179\\
8 & curvedCrossing & 111 & 0 & 33 & 1.0\% & 88 / 178\\
15 & headOn & 23 & 0 & 17 & 35.5\% & 32 / 72\\
15 & acceleratingHeadOn & 22 & 0 & 17 & 32.8\% & 46 / 104\\
15 & brakingLead & 101 & 23 & 129 & 6.1\% & 96 / 202\\
15 & crossing & 31 & 0 & 30 & 13.0\% & 44 / 100\\
15 & turningCrossing & 61 & 0 & 31 & 16.2\% & 63 / 136\\
15 & curvedHeadOn & 49 & 0 & 202 & 0.2\% & 0 / 261\\
15 & curvedCrossing & 31 & 0 & 26 & 2.4\% & 52 / 115\\
\bottomrule
\end{tabular}
\end{center}

\paragraph{Where the law violates the CLF decrease.} Head-on rollouts of the law
were started from every fifth recorded state and run for 300 steps. Grouped by $V$ at the start
of each step, no violation occurs for $V<1$ (11{,}126 steps at 8~m/s, 5{,}417 at
15~m/s). Violations occur for $1\le V<100$ (3 of 936; 22 of 284) and
$V\ge100$ (104 of 538; 203 of 299), that is, in the displaced states of an
avoidance (\path{clfViolation.m.txt}). The LQR law is a CLF law near the path,
not globally.

\paragraph{Assessment.}
\begin{enumerate}
\item \emph{Computable.} Under the target contract, membership of an endpoint
  is one deterministic rollout of the law and the target. Its linearization
  is analytic, $K\,\partial e/\partial x$, so the condensed tail rows are
  cheap.
\item \emph{Reachable without a level condition.} Requiring only that the law's
  rollout avoids the target, keeps 0.2~m and stays on the road, the median free
  horizon is 0. The maximum is 14--33 steps, except brakingLead (148, 129) and
  the 15-m/s curved head-on run (202). Under the exact models, shifting the plan
  and appending one law step keeps the endpoint in the set. That gives
  recursive feasibility of the terminal condition, and it does not depend on
  the law decreasing $V$. The CLF decrease of the law is a convergence-rate
  property, not part of the safety argument.
\item \emph{With $V\le1$ the horizon grows.} Adding the sublevel set in which
  the law met the decrease requires the ego to be back near the path at the
  endpoint. The free horizon then rises to a median of up to 105 (0 to 105) and a maximum of 72--261, longer
  than the present terminal (22--111).
\end{enumerate}
The remaining long cases are brakingLead (the law returns the ego behind the
lead) and the curved road, where the infinite-time set is empty and a 10-s
window was used. The recorded trajectories are a proxy for planned ones.

\section{Addendum 3: the CLF-constrained optimization itself as the continuation}
The user corrected Addendum~2: the design has no LQR law. The ego is driven by
optimization under the CLF constraint. The only LQR element in the code is the
quadratic matrix $P$ of the CLF, computed by \texttt{dlqr} in
\texttt{nonlinearBicycleModel.cruise}. The faithful terminal set is therefore
the set of states from which the controller itself, run without a target (CLF
constraint and road rows only), never collides with the target under its
contract. \path{scripts/studyClfOptimizationTerminal.m} simulates this
continuation by calling \texttt{collisionAvoidanceController} without a target
at every sample, from a reset, on the RK4 plant. Every second recorded
encounter state was tested. A continuation qualifies if it keeps 0.2~m from the
target, stays on the road, returns a solution at every sample and reaches the
certificate; on the curved road it must stay clear for 10~s.

\begin{center}
\begin{tabular}{rlrrrrrr}
\toprule
 & & & & separating & free & tail & check (s)\\
speed & scenario & sampled & in set & max & med/max & med/max & med/max\\
\midrule
8 & headOn & 94 & 87 & 35 & 0/14 & 64/133 & 0.19/0.43\\
8 & acceleratingHeadOn & 94 & 87 & 34 & 0/14 & 65/135 & 0.19/0.43\\
8 & brakingLead & 113 & 39 & 98 & 36/148 & 17/52 & 0.05/0.17\\
8 & crossing & 93 & 79 & 29 & 0/28 & 57/137 & 0.14/0.44\\
8 & turningCrossing & 94 & 80 & 80 & 0/28 & 107/186 & 0.30/0.60\\
8 & curvedHeadOn & 94 & 87 & 38 & 0/14 & 200/200 & 0.69/0.71\\
8 & curvedCrossing & 91 & 78 & 111 & 0/26 & 200/200 & 0.68/0.76\\
15 & headOn & 41 & 32 & 23 & 0/16 & 31/98 & 0.11/0.54\\
15 & acceleratingHeadOn & 59 & 52 & 22 & 0/14 & 31/100 & 0.14/0.56\\
15 & brakingLead & 107 & 39 & 101 & 30/136 & 23/58 & 0.10/0.32\\
15 & crossing & 57 & 44 & 31 & 0/26 & 25/106 & 0.11/0.59\\
15 & turningCrossing & 74 & 58 & 61 & 0/32 & 41/112 & 0.15/0.67\\
15 & curvedHeadOn & 599 & 488 & 49 & 0/204 & 200/200 & 1.25/1.32\\
15 & curvedCrossing & 64 & 51 & 31 & 0/26 & 200/200 & 1.21/1.24\\
\bottomrule
\end{tabular}
\end{center}
Every rejection was a predicted contact; no continuation failed to solve or
left the road. The free-horizon resolution is two samples
(\path{clf-optimization.csv}, \path{clf-optimization.log}).

\paragraph{Assessment.}
\begin{enumerate}
\item \emph{The set is useful.} Most encounter states are in it. The free horizon
  needed to reach it is at most 14--32 samples, against 22--111 for the
  separating terminal, except for brakingLead (148, 136) and the 15-m/s curved
  head-on run (204). These are close to the values of Addendum~2, so the LQR
  proxy described the optimization controller well where it qualified.
\item \emph{It cannot be evaluated online by simulation.} One membership check
  of a single endpoint takes 0.05--1.3~s, because the continuation is a
  sequence of up to 200 conic solves. An optimization needs more than membership
  of one point: it needs a description of the set near the endpoint (rows or a
  linearization). The closed-loop map of a sequence of conic programs is
  piecewise and has no closed form. Its sensitivity would require propagating
  KKT sensitivities through every solve.
\item \emph{It is target independent.} After a reset the continuation of the ego
  depends only on its state relative to the path (lateral error, heading
  error, speeds, yaw rate and the held braking ratio, with free path phase).
  It does not depend on the target. It can therefore be computed offline, per
  road curvature, as a table or fitted model of the ego's continuation in path
  coordinates. Online, membership and its endpoint sensitivity would then be a
  table evaluation composed with the closed-form target prediction. This was
  not built; table size, interpolation error and the dependence on the
  controller's warm-start memory are open questions.
\end{enumerate}
\end{document}
